We presented ComplianceAI, a retrieval-augmented generation system for
automatically checking whether university policy documents comply with YÖK
regulations. The system was motivated by a genuine administrative need:
compliance review at Turkish universities is a manual, time-consuming process
that requires reading through dense legal text.
We built something that can meaningfully assist administrative staff ---
not replace them, but make their work faster and more targeted.

\subsection{Summary of Findings}

Our evaluation on a 40-case RAG-aware benchmark produced clear,
reproducible findings.
Retrieval helps: the best RAG configurations reach \textbf{47.5\% accuracy}
compared to \textbf{37.5\% for no-RAG}, a 10 percentage point improvement
on a benchmark specifically designed to require regulatory knowledge
(see Table~\ref{tab:main_results} and Figure~\ref{fig:accuracy}).
Sparse and hybrid retrieval outperform dense-only retrieval in this domain,
likely because BM25's exact-match scoring handles Turkish legal terminology
well without relying on embedding quality.
GPT-4o-mini matches GPT-4o performance at 19$\times$ lower cost ---
a significant practical advantage for a system meant to run continuously
in a university environment.
Optimal hyperparameters are 350-word chunks and Top-5 retrieval depth,
as confirmed by the ablation studies in
Tables~\ref{tab:chunk_ablation} and \ref{tab:topk_ablation} and
Figures~\ref{fig:chunk_ablation} and \ref{fig:topk_ablation}.

The most important limitation is that all models systematically fail to
detect non-compliant documents, defaulting to ``partial'' even for clear
violations (Table~\ref{tab:perclass}, Figure~\ref{fig:heatmap}).
This means the system is currently best positioned as a \textit{risk filter}
that flags documents for human review, not as a standalone compliance auditor.

\subsection{Limitations}

\textbf{Benchmark size.}
Our test set contains 40 cases. While this is sufficient to observe
consistent trends, a larger annotated dataset --- ideally with human expert
labels rather than LLM-generated gold labels --- would give more reliable
performance estimates.

\textbf{Single institution.}
We built and evaluated the system using Istanbul Bilgi University documents.
Whether the findings generalize across institutions with different policy
styles and document formats is an open question.

\textbf{Manual document management.}
The current knowledge base requires manual updates whenever YÖK revises
a regulation. This is manageable for a prototype but not scalable
for a production system used across many institutions.

\textbf{Non-compliant detection.}
Models consistently under-predict the non-compliant class.
This is a fundamental behavior of current LLMs, not a hyperparameter problem,
and will likely require fine-tuning or a different problem formulation to fix.

\subsection{Future Work}

Several directions are worth pursuing.

\paragraph{Live regulatory updates.}
The most immediate improvement would be to automate regulatory document
updates by integrating directly with the YÖK web portal.
When YÖK publishes a new or revised regulation, the system would
automatically fetch the document, chunk it, update the indexes, and notify
administrators. This would eliminate the current manual maintenance step
and ensure the system always operates on up-to-date regulations.

\paragraph{Broader deployment.}
Although ComplianceAI was built for Istanbul Bilgi University, it is designed
to be institution-agnostic. Any Turkish university could deploy the same system
by loading their own institutional documents alongside the shared YÖK knowledge
base. Beyond Turkey, the same approach is directly applicable to other
higher education regulatory systems --- ECTS credit regulations in Europe,
ABET accreditation requirements in the United States, or country-specific
higher education laws elsewhere.

\paragraph{Better non-compliant detection.}
Fine-tuning a smaller language model on compliance-specific examples, or
using adversarial training with deliberately non-compliant document pairs,
could improve recall for the non-compliant class.

\paragraph{Graph-based retrieval.}
YÖK regulations frequently cross-reference each other.
Graph-based retrieval approaches, as explored by \citet{Masoudifard2024},
could better capture these relationships and improve performance on
multi-article compliance questions.

\paragraph{Human annotation.}
Constructing a human-annotated benchmark with cases drawn from real
administrative documents would significantly strengthen evaluation reliability
and provide a community resource for future work in this area.

\subsection{Closing Remarks}

ComplianceAI demonstrates that retrieval-augmented generation is a practical
and effective approach for regulatory compliance checking in a specialized
legal domain. The system is not yet production-ready as a standalone tool,
but it meaningfully accelerates the compliance review process by surfacing
relevant regulation passages and providing an initial classification.
We hope the open-source code, benchmark dataset, and evaluation methodology
we provide will serve as a foundation for future work --- whether at Turkish
universities seeking to modernize their administrative processes,
or in other regulatory domains where the same challenges apply.
